\documentclass[11pt,a4paper]{article}
% preamble.tex — shared setup for human-verification documents
% Compile each document with:  tectonic <File>.tex

\usepackage[T1]{fontenc}
\usepackage[utf8]{inputenc}
\usepackage{lmodern}
\usepackage{amssymb,amsmath,amsthm}
\usepackage{mathtools}
\usepackage[dvipsnames]{xcolor}
\usepackage{listings}
\usepackage{hyperref}
\usepackage[margin=2.5cm]{geometry}
\usepackage{parskip}

% ── Colours ──────────────────────────────────────────────────────────
\definecolor{leanBg}{HTML}{F7F7F0}
\definecolor{haskellBg}{HTML}{F0F4F7}
\definecolor{keyword}{HTML}{7B3294}
\definecolor{comment}{HTML}{5A8C32}
\definecolor{noteFg}{HTML}{B35900}

% ── Lean 4 listings language ─────────────────────────────────────────
% NOTE: Lean listing blocks use ASCII approximations of unicode symbols
% (e.g. -> for arrow, o. for compose).  The literate key substitutes
% these with proper math-mode renderings in the PDF.
\lstdefinelanguage{Lean4}{
  morekeywords={def,theorem,structure,where,namespace,end,import,fun,
                let,have,show,by,exact,intro,cases,rfl,subst,obtain,
                funext,simp,rw,apply,variable,Type,Prop,instance,
                section,open,class,extends,deriving,sorry,noncomputable},
  sensitive=true,
  morecomment=[l]{--},
  morecomment=[n]{/-}{-/},
  morestring=[b]",
}

% ── Lean listing style ───────────────────────────────────────────────
\lstdefinestyle{lean}{
  language=Lean4,
  backgroundcolor=\color{leanBg},
  basicstyle=\small\ttfamily,
  keywordstyle=\color{keyword}\bfseries,
  commentstyle=\color{comment}\itshape,
  breaklines=true,
  frame=single,
  framerule=0.4pt,
  xleftmargin=4pt,
  xrightmargin=4pt,
  aboveskip=6pt,
  belowskip=4pt,
  literate=
    {->}{{$\to$}}2
    {=>}{{$\Rightarrow$}}2
    {<|}{{$\langle$}}1
    {|>}{{$\rangle$}}1
    {==>}{{$\Longrightarrow$}}3
    { o. }{{$\;\circ\;$}}3
    {forall}{{$\forall$}}1
    {:=}{{$\,:=\,$}}2
    {!=}{{$\neq$}}2
    {<=}{{$\leq$}}2
    {>=}{{$\geq$}}2,
}

% ── Haskell listing style ───────────────────────────────────────────
\lstdefinestyle{haskell}{
  language=Haskell,
  backgroundcolor=\color{haskellBg},
  basicstyle=\small\ttfamily,
  keywordstyle=\color{keyword}\bfseries,
  commentstyle=\color{comment}\itshape,
  breaklines=true,
  frame=single,
  framerule=0.4pt,
  xleftmargin=4pt,
  xrightmargin=4pt,
  aboveskip=6pt,
  belowskip=4pt,
  literate=
    {->}{{$\to$}}2
    {=>}{{$\Rightarrow$}}2,
}

% ── Theorem environments ────────────────────────────────────────────
\theoremstyle{definition}
\newtheorem{defn}{Definition}[section]
\newtheorem{thm}[defn]{Theorem}
\newtheorem{prop}[defn]{Proposition}

\theoremstyle{remark}
\newtheorem*{fidelitynote}{Fidelity note}

% ── Convenience macros ───────────────────────────────────────────────
\newcommand{\lean}[1]{\texttt{#1}}      % inline Lean identifier
\newcommand{\hask}[1]{\texttt{#1}}      % inline Haskell identifier
\newcommand{\Ext}{\mathrm{Ext}}
\newcommand{\Pos}{\mathrm{Pos}}
\newcommand{\Shape}{\mathrm{Shape}}


\title{Human Verification: \texttt{Containers.Monoidal}\\[4pt]
  \large The Sequential Operator as a Monoidal Category}
\author{Machine-generated from \texttt{lean/Monoidal.lean} for review by Neil Ghani}
\date{July 2026}

\begin{document}
\maketitle

\begin{abstract}
This document presents every definition and theorem \emph{statement} from the
Lean~4 module \lean{Containers.Monoidal} in three parallel forms: the original
Lean, a readable Haskell rendering, and precise mathematical English.  Proofs
are omitted --- they are machine-checked by Lean.  The purpose is to let a
human verifier confirm that the formalisation captures the intended
mathematics.

The file proves that $(\mathbf{Cont}, \triangleleft, I)$ is a monoidal
category, by supplying the bifunctoriality of $\triangleleft$ on morphisms,
the naturality of the structural isomorphisms, the pentagon, and the triangle.
All laws hold by \lean{rfl} (definitional equality), since the position fibres
line up definitionally.  A minimal Mathlib-free \lean{MonoidalCategory}
typeclass is also defined, and the instance \lean{ContMonoidal} assembles the
whole.
\end{abstract}

\tableofcontents
\bigskip

%===================================================================
\section{Overview}
%===================================================================

Building on \texttt{Containers.Sequential}, which defines $\triangleleft$, the
unit $I$, and the associator/unitors as isomorphisms, this file supplies the
remaining ingredients to certify $(\mathbf{Cont}, \triangleleft, I)$ as a
monoidal category.

Key insight: the associator's shape map is \emph{currying} and its position map
is $\Sigma$-reassociation.  Both pentagon legs reduce to the same flattening of
a fourfold composite, and every map is the identity on the underlying flat
tuple.

%===================================================================
\section{Bifunctoriality of \texorpdfstring{$\triangleleft$}{seq} on morphisms}
%===================================================================

\subsection{Definition: \lean{Container.seq\textsubscript{2}}}

\begin{defn}[Bifunctorial action of $\triangleleft$]
Given morphisms $\mu : G \Rightarrow G'$ and $\nu : F \Rightarrow F'$, the
\textbf{action of $\triangleleft$ on the pair} is the morphism
$\lean{seq}_2\,\mu\,\nu : G \triangleleft F \Rightarrow G' \triangleleft F'$.
On shapes, a shape $(t, f)$ of $G \triangleleft F$ --- an outer shape $t$ and
an inner-shape assignment $f : P_G(t) \to S_F$ --- maps to
$(\mu(t),\; \nu \circ f \circ \mu^\sharp_t)$, where $\mu^\sharp_t$ is the
backward position map of $\mu$ at $t$.
On positions: both backward position maps are applied componentwise.
\end{defn}

\paragraph{Lean.}
\begin{lstlisting}[style=lean]
def Container.seq2 {G G' F F' : Container}
    (mu : ContainerMorphism G G') (nu : ContainerMorphism F F') :
    ContainerMorphism (G seq F) (G' seq F') where
  onShapes := fun tf =>
    <|mu.onShapes tf.1, fun q' => nu.onShapes (tf.2 (mu.onPos tf.1 q'))|>
  onPos := fun tf qr =>
    <|mu.onPos tf.1 qr.1, nu.onPos (tf.2 (mu.onPos tf.1 qr.1)) qr.2|>
\end{lstlisting}

\paragraph{Haskell.}
\begin{lstlisting}[style=haskell]
-- G seq F has Shape = (t, f) where t :: Shape G, f :: Pos G t -> Shape F
-- seq2 mu nu : ContainerMorphism (G `seq` F) (G' `seq` F')
seq2 :: Morph g p g' p' -> Morph f q f' q'
     -> Morph (g `Seq` f) (p, q) (g' `Seq` f') (p', q')
seq2 mu nu = Morph
  { onShapes = \(t, f) -> (onShapes mu t, \q' -> onShapes nu (f (onPos mu t q')))
  , onPos    = \(t, f) (qr1, qr2) ->
      (onPos mu t qr1, onPos nu (f (onPos mu t qr1)) qr2) }
\end{lstlisting}

\begin{fidelitynote}
The Lean definition uses dependent types: \lean{tf.2} is a function from
$P_G(\lean{tf.1})$, whose type depends on the shape $\lean{tf.1}$.  The Haskell
rendering non-dependently pairs the components.
\end{fidelitynote}


\subsection{Definition: \lean{Container.whiskerLeft}}

\begin{defn}[Left whiskering]
\textbf{Left whiskering} $G \triangleleft \nu : G \triangleleft F \Rightarrow G \triangleleft F'$
acts by $\nu$ on the inner container, leaving the outer container $G$ fixed:
\[
  G \triangleleft \nu \;=\; \lean{seq}_2\,(\mathrm{id}_G)\,\nu.
\]
\end{defn}

\paragraph{Lean.}
\begin{lstlisting}[style=lean]
def Container.whiskerLeft (G : Container) {F F' : Container}
    (nu : ContainerMorphism F F') : ContainerMorphism (G seq F) (G seq F') :=
  Container.seq2 (ContainerMorphism.id G) nu
\end{lstlisting}

\paragraph{Haskell.}
\begin{lstlisting}[style=haskell]
whiskerLeft :: Proxy g -> Morph f q f' q'
            -> Morph (g `Seq` f) (p, q) (g `Seq` f') (p, q')
whiskerLeft _ nu = seq2 idMorph nu
\end{lstlisting}


\subsection{Definition: \lean{Container.whiskerRight}}

\begin{defn}[Right whiskering]
\textbf{Right whiskering} $\mu \triangleleft F : G \triangleleft F \Rightarrow G' \triangleleft F$
acts by $\mu$ on the outer container, leaving the inner container $F$ fixed:
\[
  \mu \triangleleft F \;=\; \lean{seq}_2\,\mu\,(\mathrm{id}_F).
\]
\end{defn}

\paragraph{Lean.}
\begin{lstlisting}[style=lean]
def Container.whiskerRight {G G' : Container} (mu : ContainerMorphism G G')
    (F : Container) : ContainerMorphism (G seq F) (G' seq F) :=
  Container.seq2 mu (ContainerMorphism.id F)
\end{lstlisting}

\paragraph{Haskell.}
\begin{lstlisting}[style=haskell]
whiskerRight :: Morph g p g' p' -> Proxy f
             -> Morph (g `Seq` f) (p, q) (g' `Seq` f) (p', q)
whiskerRight mu _ = seq2 mu idMorph
\end{lstlisting}


%===================================================================
\section{Bifunctor laws for \texorpdfstring{$\triangleleft$}{seq}}
%===================================================================

\subsection{Theorem: \lean{Container.seq\textsubscript{2}\_id}}

\begin{thm}[Bifunctoriality: identity]
$\lean{seq}_2$ preserves identities:
\[
  \lean{seq}_2\,(\mathrm{id}_G)\,(\mathrm{id}_F) \;=\; \mathrm{id}_{G \triangleleft F}.
\]
\end{thm}

\paragraph{Lean.}
\begin{lstlisting}[style=lean]
theorem Container.seq2_id (G F : Container) :
    Container.seq2 (ContainerMorphism.id G) (ContainerMorphism.id F)
      = ContainerMorphism.id (G seq F) := rfl
\end{lstlisting}

\paragraph{Haskell.}
\begin{lstlisting}[style=haskell]
-- seq2 idMorph idMorph = idMorph
\end{lstlisting}


\subsection{Theorem: \lean{Container.seq\textsubscript{2}\_comp}}

\begin{thm}[Bifunctoriality: interchange]
$\lean{seq}_2$ preserves composition (interchange / bifunctoriality):
\[
  \lean{seq}_2\,(\mu \cdot \mu')\,(\nu \cdot \nu')
  \;=\;
  \lean{seq}_2\,\mu\,\nu \;\cdot\; \lean{seq}_2\,\mu'\,\nu'.
\]
\end{thm}

\paragraph{Lean.}
\begin{lstlisting}[style=lean]
theorem Container.seq2_comp {G G' G'' F F' F'' : Container}
    (mu : ContainerMorphism G G') (mu' : ContainerMorphism G' G'')
    (nu : ContainerMorphism F F') (nu' : ContainerMorphism F' F'') :
    Container.seq2 (mu.comp mu') (nu.comp nu')
      = (Container.seq2 mu nu).comp (Container.seq2 mu' nu') := rfl
\end{lstlisting}

\paragraph{Haskell.}
\begin{lstlisting}[style=haskell]
-- seq2 (mu `compM` mu') (nu `compM` nu')
--   = seq2 mu nu `compM` seq2 mu' nu'
\end{lstlisting}


%===================================================================
\section{Naturality of the structural isomorphisms}
%===================================================================

\subsection{Theorem: \lean{Container.associator\_naturality}}

\begin{thm}[Naturality of the associator]
The associator $\alpha$ is natural in all three arguments:
for all morphisms $f_i : X_i \Rightarrow Y_i$,
\[
  (\lean{seq}_2\,(\lean{seq}_2\,f_1\,f_2)\,f_3) \cdot \alpha_{Y_1,Y_2,Y_3}
  \;=\;
  \alpha_{X_1,X_2,X_3} \cdot \lean{seq}_2\,f_1\,(\lean{seq}_2\,f_2\,f_3).
\]
\end{thm}

\paragraph{Lean.}
\begin{lstlisting}[style=lean]
theorem Container.associator_naturality {X1 Y1 X2 Y2 X3 Y3 : Container}
    (f1 : ContainerMorphism X1 Y1) (f2 : ContainerMorphism X2 Y2)
    (f3 : ContainerMorphism X3 Y3) :
    (Container.seq2 (Container.seq2 f1 f2) f3).comp
        (Container.associator Y1 Y2 Y3).hom
      = (Container.associator X1 X2 X3).hom.comp
          (Container.seq2 f1 (Container.seq2 f2 f3)) :=
  rfl
\end{lstlisting}

\paragraph{Haskell.}
\begin{lstlisting}[style=haskell]
-- (seq2 (seq2 f1 f2) f3) `compM` assocHom y1 y2 y3
--   = assocHom x1 x2 x3 `compM` seq2 f1 (seq2 f2 f3)
\end{lstlisting}


\subsection{Theorem: \lean{Container.leftUnitor\_naturality}}

\begin{thm}[Naturality of the left unitor]
The left unitor $\lambda$ is natural: for $f : X \Rightarrow Y$,
\[
  (I \triangleleft f) \cdot \lambda_Y \;=\; \lambda_X \cdot f.
\]
\end{thm}

\paragraph{Lean.}
\begin{lstlisting}[style=lean]
theorem Container.leftUnitor_naturality {X Y : Container}
    (f : ContainerMorphism X Y) :
    (Container.whiskerLeft Container.I f).comp (Container.leftUnitor Y).hom
      = (Container.leftUnitor X).hom.comp f := rfl
\end{lstlisting}

\paragraph{Haskell.}
\begin{lstlisting}[style=haskell]
-- whiskerLeft I f `compM` leftUnitorHom y
--   = leftUnitorHom x `compM` f
\end{lstlisting}


\subsection{Theorem: \lean{Container.rightUnitor\_naturality}}

\begin{thm}[Naturality of the right unitor]
The right unitor $\rho$ is natural: for $f : X \Rightarrow Y$,
\[
  (f \triangleleft I) \cdot \rho_Y \;=\; \rho_X \cdot f.
\]
\end{thm}

\paragraph{Lean.}
\begin{lstlisting}[style=lean]
theorem Container.rightUnitor_naturality {X Y : Container}
    (f : ContainerMorphism X Y) :
    (Container.whiskerRight f Container.I).comp (Container.rightUnitor Y).hom
      = (Container.rightUnitor X).hom.comp f := rfl
\end{lstlisting}

\paragraph{Haskell.}
\begin{lstlisting}[style=haskell]
-- whiskerRight f I `compM` rightUnitorHom y
--   = rightUnitorHom x `compM` f
\end{lstlisting}


%===================================================================
\section{Coherence}
%===================================================================

\subsection{Theorem: \lean{Container.pentagon}}

\begin{thm}[Pentagon identity]
The two ways of re-associating
$((C_1 \triangleleft C_2) \triangleleft C_3) \triangleleft C_4$
to $C_1 \triangleleft (C_2 \triangleleft (C_3 \triangleleft C_4))$ coincide:
\[
  (\alpha_{C_1,C_2,C_3} \triangleleft C_4) \cdot
  \alpha_{C_1, C_2 \triangleleft C_3, C_4} \cdot
  (C_1 \triangleleft \alpha_{C_2,C_3,C_4})
  \;=\;
  \alpha_{C_1 \triangleleft C_2, C_3, C_4} \cdot
  \alpha_{C_1, C_2, C_3 \triangleleft C_4}.
\]
\end{thm}

\paragraph{Lean.}
\begin{lstlisting}[style=lean]
theorem Container.pentagon (C1 C2 C3 C4 : Container) :
    ((Container.whiskerRight (Container.associator C1 C2 C3).hom C4).comp
        (Container.associator C1 (C2 seq C3) C4).hom).comp
        (Container.whiskerLeft C1 (Container.associator C2 C3 C4).hom)
      = (Container.associator (C1 seq C2) C3 C4).hom.comp
        (Container.associator C1 C2 (C3 seq C4)).hom := rfl
\end{lstlisting}

\paragraph{Haskell.}
\begin{lstlisting}[style=haskell]
-- (whiskerRight assoc123 C4 `compM` assoc1_23_4 `compM`
--   whiskerLeft C1 assoc234)
--   = assoc12_3_4 `compM` assoc1_2_34
\end{lstlisting}


\subsection{Theorem: \lean{Container.triangle}}

\begin{thm}[Triangle identity]
The associator through the unit composed with the inner left unitor equals
the outer right unitor:
\[
  \alpha_{C_1, I, C_2} \cdot (C_1 \triangleleft \lambda_{C_2})
  \;=\;
  \rho_{C_1} \triangleleft C_2,
\]
as morphisms $(C_1 \triangleleft I) \triangleleft C_2 \Rightarrow C_1 \triangleleft C_2$.
\end{thm}

\paragraph{Lean.}
\begin{lstlisting}[style=lean]
theorem Container.triangle (C1 C2 : Container) :
    (Container.associator C1 Container.I C2).hom.comp
        (Container.whiskerLeft C1 (Container.leftUnitor C2).hom)
      = Container.whiskerRight (Container.rightUnitor C1).hom C2 := rfl
\end{lstlisting}

\paragraph{Haskell.}
\begin{lstlisting}[style=haskell]
-- assocHom c1 I c2 `compM` whiskerLeft c1 (leftUnitorHom c2)
--   = whiskerRight (rightUnitorHom c1) c2
\end{lstlisting}


%===================================================================
\section{The monoidal category typeclass}
%===================================================================

\subsection{Structure: \lean{CatIso}}

\begin{defn}[Isomorphism in a category]
A \textbf{categorical isomorphism} between objects $X$ and $Y$ in a category
$\mathcal{C}$ consists of a forward morphism \lean{hom} and a backward morphism
\lean{inv} satisfying both round-trip conditions:
$\lean{hom} \cdot \lean{inv} = \mathrm{id}_X$ and
$\lean{inv} \cdot \lean{hom} = \mathrm{id}_Y$.
\end{defn}

\paragraph{Lean.}
\begin{lstlisting}[style=lean]
structure CatIso {Obj : Type u} [Category Obj] (X Y : Obj) where
  hom : Category.Hom X Y
  inv : Category.Hom Y X
  hom_inv_id : Category.comp hom inv = Category.id X
  inv_hom_id : Category.comp inv hom = Category.id Y
\end{lstlisting}

\paragraph{Haskell.}
\begin{lstlisting}[style=haskell]
data CatIso cat x y = CatIso
  { catHom    :: cat x y
  , catInv    :: cat y x
  , homInvId  :: catHom `comp` catInv == idC x  -- propositional
  , invHomId  :: catInv `comp` catHom == idC y
  }
\end{lstlisting}


\subsection{Class: \lean{MonoidalCategory}}

\begin{defn}[Monoidal category (Mathlib-free)]
A \textbf{monoidal category} extends \lean{Category} with:
\begin{itemize}
  \item a tensor on objects $\otimes : \mathrm{Obj} \to \mathrm{Obj} \to \mathrm{Obj}$,
  \item a bifunctorial tensor on morphisms,
  \item a unit object $\mathbf{1}$,
  \item associator and unitor isomorphisms,
  \item naturality of the associator and unitors, and
  \item the pentagon and triangle coherence laws.
\end{itemize}
\end{defn}

\paragraph{Lean.}
\begin{lstlisting}[style=lean]
class MonoidalCategory (Obj : Type u) extends Category.{u, v} Obj where
  tensorObj    : Obj -> Obj -> Obj
  tensorHom    : {X1 Y1 X2 Y2 : Obj} ->
    Hom X1 Y1 -> Hom X2 Y2 -> Hom (tensorObj X1 X2) (tensorObj Y1 Y2)
  tensorHom_id : forall (X Y : Obj),
    tensorHom (id X) (id Y) = id (tensorObj X Y)
  tensorHom_comp : forall {X1 Y1 Z1 X2 Y2 Z2 : Obj}
    (f1 : Hom X1 Y1) (g1 : Hom Y1 Z1) (f2 : Hom X2 Y2) (g2 : Hom Y2 Z2),
    tensorHom (comp f1 g1) (comp f2 g2)
      = comp (tensorHom f1 f2) (tensorHom g1 g2)
  tensorUnit       : Obj
  associator       : (X Y Z : Obj) ->
    CatIso (tensorObj (tensorObj X Y) Z) (tensorObj X (tensorObj Y Z))
  leftUnitor       : (X : Obj) -> CatIso (tensorObj tensorUnit X) X
  rightUnitor      : (X : Obj) -> CatIso (tensorObj X tensorUnit) X
  associator_naturality : forall {X1 Y1 X2 Y2 X3 Y3 : Obj}
    (f1 : Hom X1 Y1) (f2 : Hom X2 Y2) (f3 : Hom X3 Y3),
    comp (tensorHom (tensorHom f1 f2) f3) (associator Y1 Y2 Y3).hom
      = comp (associator X1 X2 X3).hom (tensorHom f1 (tensorHom f2 f3))
  leftUnitor_naturality : forall {X Y : Obj} (f : Hom X Y),
    comp (tensorHom (id tensorUnit) f) (leftUnitor Y).hom
      = comp (leftUnitor X).hom f
  rightUnitor_naturality : forall {X Y : Obj} (f : Hom X Y),
    comp (tensorHom f (id tensorUnit)) (rightUnitor Y).hom
      = comp (rightUnitor X).hom f
  pentagon : forall (W X Y Z : Obj),
    comp (comp (tensorHom (associator W X Y).hom (id Z))
               (associator W (tensorObj X Y) Z).hom)
         (tensorHom (id W) (associator X Y Z).hom)
      = comp (associator (tensorObj W X) Y Z).hom
             (associator W X (tensorObj Y Z)).hom
  triangle : forall (X Y : Obj),
    comp (associator X tensorUnit Y).hom
         (tensorHom (id X) (leftUnitor Y).hom)
      = tensorHom (rightUnitor X).hom (id Y)
\end{lstlisting}

\paragraph{Haskell.}
\begin{lstlisting}[style=haskell]
class Category cat => MonoidalCategory cat where
  tensorObj  :: obj -> obj -> obj
  tensorHom  :: cat x1 y1 -> cat x2 y2
             -> cat (tensorObj x1 x2) (tensorObj y1 y2)
  tensorUnit :: obj
  assoc      :: CatIso cat (tensorObj (tensorObj x y) z)
                            (tensorObj x (tensorObj y z))
  leftUnit   :: CatIso cat (tensorObj tensorUnit x) x
  rightUnit  :: CatIso cat (tensorObj x tensorUnit) x
  -- plus naturality, pentagon, triangle laws as proof obligations
\end{lstlisting}


\subsection{Definition: \lean{ContainerIso.toCatIso}}

\begin{defn}[Reinterpretation of container isomorphisms]
A \lean{ContainerIso} (from \texttt{Cont.lean}) can be viewed as a \lean{CatIso}
for the category \lean{Container}: the fields coincide definitionally.
\end{defn}

\paragraph{Lean.}
\begin{lstlisting}[style=lean]
def ContainerIso.toCatIso {C D : Container} (i : ContainerIso C D) :
    CatIso C D :=
  <|i.hom, i.inv, i.hom_inv, i.inv_hom|>
\end{lstlisting}

\paragraph{Haskell.}
\begin{lstlisting}[style=haskell]
toCatIso :: ContainerIso c d -> CatIso ContainerMorph c d
toCatIso i = CatIso (isoHom i) (isoInv i) (homInv i) (invHom i)
\end{lstlisting}


\subsection{Instance: \lean{ContMonoidal}}

\begin{defn}[$({\bf Cont}, \triangleleft, I)$ is monoidal]
The sequential operator $\triangleleft$ with morphism action $\lean{seq}_2$, unit
container $I$, and structural isomorphisms from \texttt{Sequential.lean} constitute
a monoidal category.  Every field --- the two functor laws, three naturality
squares, pentagon, triangle --- is discharged by \lean{rfl}.
\end{defn}

\paragraph{Lean.}
\begin{lstlisting}[style=lean]
instance ContMonoidal : MonoidalCategory Container where
  tensorObj         := Container.seq
  tensorHom         := Container.seq2
  tensorHom_id      := Container.seq2_id
  tensorHom_comp    := fun f1 g1 f2 g2 => Container.seq2_comp f1 g1 f2 g2
  tensorUnit        := Container.I
  associator        := fun X Y Z => (Container.associator X Y Z).toCatIso
  leftUnitor        := fun X => (Container.leftUnitor X).toCatIso
  rightUnitor       := fun X => (Container.rightUnitor X).toCatIso
  associator_naturality :=
    fun f1 f2 f3 => Container.associator_naturality f1 f2 f3
  leftUnitor_naturality  := fun f => Container.leftUnitor_naturality f
  rightUnitor_naturality := fun f => Container.rightUnitor_naturality f
  pentagon          := Container.pentagon
  triangle          := Container.triangle
\end{lstlisting}

\paragraph{Haskell.}
\begin{lstlisting}[style=haskell]
-- The sequential operator `seq` (written `triangleleft`) gives Cont
-- the structure of a monoidal category, with:
--   tensorObj = seq
--   tensorHom = seq2
--   unit      = I
-- All laws hold definitionally.
instance MonoidalCategory ContainerMorph where ...
\end{lstlisting}

\begin{fidelitynote}
In Lean, only one \lean{MonoidalCategory Container} \emph{instance} can be
active at a time; this is the canonical one.  The other three monoidal
structures on \lean{Container} (coproduct $+$, product $\times$, Dirichlet $\otimes$)
are given as named \lean{def}s in \texttt{FourMonoidal.lean} /
\texttt{Dirichlet.lean}, to avoid instance clashes.
\end{fidelitynote}

\end{document}
